% ============ 第 2 讲 贸易模型导论与禀赋模型 ============
\chap{第 2 讲　贸易模型导论与禀赋模型}{Introduction to Trade Models \& The Endowment Model}

\textbf{本讲导读}：本讲对应 Lecture 2（Introduction to Trade Models，Sep.\ 11），共 61 页，分两大部分：
第一部分给出\textbf{所有古典贸易模型共享的设定}（两个国家、两种商品、供给面、需求面）与求解贸易均衡的\textbf{关键条件（相对价格均等化）}；
第二部分在最简单的环境（\textbf{无生产、纯禀赋经济}）中完整演示贸易模型的结构——从自给自足均衡、套利、贸易均衡、贸易模式、贸易平衡到贸易得益，并一般化到 $G$ 种商品、$N$ 个国家，应用于石油贸易。课后笔记页（第 54--61 页，手写）给出了极佳的模型结构总结，本章``方法地图''即以它为骨架。

\secn{2.1 为什么要建模：模型的对象与结构}

\textbf{为什么经济学要用模型？（第 3 页）}
\begin{itemize}
  \item 对现实的\term{简化}（simplification，即假设）；
  \item 分析复杂的逻辑关系；
  \item \kw{可检验的预测（testable predictions）}（荧光高亮；很多情况下是反直觉的）。
\end{itemize}

\textbf{我们建模的对象是什么？（第 4 页）}
\begin{itemize}
  \item \kw{设定/情境（Setting/context）}：决策者是谁？\textbf{选择集}是什么？（老师划线强调）交易条件如何？
  \item 决策者的\term{目标与约束}（objectives and constraints）；
  \item 更复杂的结构：信息？交易成本？不确定性？策略互动？——老师手写批注``\kw{Geography}''（地理也是这类``背景结构''的一种，第 6 页标题旁再次手写``Geography''）。
\end{itemize}

\begin{exbox}[板书示例：厂商问题中的内生与外生变量（第 4 页手写）]
考察厂商 $i$：
\begin{itemize}
  \item \kw{内生变量（endogenous）}：$L_i$（劳动）、$K_i$（资本）、$Q_i$（产量）、$P_i \in [0,+\infty)$（价格）；
  \item \term{外生变量（exogenous）}：$w_i$（工资）、$r_i$（利率）、$F(K,L)$（生产技术）、$P(Q_i)$（反需求）。
\end{itemize}
两步优化：
\[
\min_{K_i, L_i}\; w_i L_i + r_i K_i \quad \text{s.t.}\quad F(K_i, L_i) = Q_i
\;\Longleftrightarrow\; C_i(Q_i) \;(\text{成本函数}),
\qquad
\max_{Q_i}\; \pi_i = P(Q_i)\, Q_i - C_i(Q_i).
\]
即先由成本最小化导出成本函数 $C_i(Q_i)$，再做利润最大化。
\end{exbox}

\begin{figure}[H]
  \centering
  \includegraphics[width=0.86\linewidth]{figs/ch02_fig_01.png}
  \caption{第 4 页板书：厂商问题中的内生/外生变量与两步优化（成本最小化 $\Rightarrow$ 成本函数 $\Rightarrow$ 利润最大化）。}
  \label{fig:ch02-firm}
\end{figure}

\secn{2.2 古典贸易模型的共同设定}

\begin{assumpbox}[共同设定（Common Setup，第 6--9 页）]
\begin{enumerate}
  \item \term{两个（或多个）国家}：至少需要两个国家才能贸易；若只对某一特定国家感兴趣，可以把另一国建模为``\term{世界其他地区（Rest of the World}''，老师手写缩写 \kw{ROW}）。
  \item \term{两种（或多种）商品}：只有一种商品时均衡无法平衡——买一样东西，就必须卖另一样来平衡预算。
  \item \term{商品供给面（supply side）}（老师板书 PPF 草图，见图~\ref{fig:ch02-ppf}）：
    \begin{itemize}
      \item \kw{禀赋（endowment）}（要素）+ \term{生产函数（技术）} $\Rightarrow$ \term{生产可能性边界（PPF）}；
      \item \kw{一个}假想生产者进行\kw{利润最大化}（荧光高亮）；
      \item 从\term{完全竞争}条件出发：\kw{MC=P，w=MRPL，r=MRPK}（荧光高亮；边际成本定价、要素价格等于边际产品价值）。
    \end{itemize}
  \item \term{商品需求面（demand side）}（老师板书预算线+无差异曲线草图，见图~\ref{fig:ch02-ppf}右）：
    \begin{itemize}
      \item \kw{一个}假想消费者代表整个国家；
      \item ``\kw{代表性主体（representative agent）}''（荧光高亮）——老师手写其成立条件：\term{Identical（相同偏好）、Homothetic（位似偏好）}时，$N$ 个消费者 $\Longleftrightarrow$ 1 个代表性主体；
      \item 对两种（或多种）商品的\kw{偏好}——\term{无差异曲线}（荧光高亮）；
      \item \kw{预算约束}（荧光高亮）。
    \end{itemize}
\end{enumerate}
\end{assumpbox}

\begin{figure}[H]
  \centering
  \begin{minipage}{0.45\linewidth}\centering
    \includegraphics[width=\linewidth]{figs/ch02_fig_02.png}
  \end{minipage}\hfill
  \begin{minipage}{0.45\linewidth}\centering
    \includegraphics[width=\linewidth]{figs/ch02_fig_03.png}
  \end{minipage}
  \caption{左：供给面板书（弯向原点的 PPF）；右：需求面板书（预算线与无差异曲线相切）。}
  \label{fig:ch02-ppf}
\end{figure}

\textbf{封闭经济（autarky，自给自足）与自由贸易（第 9 页）}：
\begin{itemize}
  \item \kw{封闭国家（autarky，老师手写``自给自足''）}：
    \begin{itemize}
      \item 是比较的\term{基准（baseline）}；
      \item \kw{预算约束就是 PPF}——你消费的就是你生产的（老师洋红下划线强调）；板书图中生产点 $(Q_A, Q_B)$ 与消费点 $(C_A, C_B)$ 重合（见图~\ref{fig:ch02-autarky}）。
    \end{itemize}
  \item \term{两国\kw{彼此自由贸易}（荧光高亮）}：
    \begin{itemize}
      \item \term{相对价格均等化}——无套利条件（no arbitrage condition）；
      \item 预算约束由相对价格决定（老师手写：$I = P_A \cdot A + P_B \cdot B$）；
      \item 福利比较——代表性主体的效用差距。
    \end{itemize}
\end{itemize}

\begin{figure}[H]
  \centering
  \includegraphics[width=0.4\linewidth]{figs/ch02_fig_04.png}
  \caption{第 9 页板书：自给自足均衡——PPF 与无差异曲线相切，生产点 $(Q_A,Q_B)$ = 消费点 $(C_A,C_B)$；预算约束（价格线）$I = P_A A + P_B B$。}
  \label{fig:ch02-autarky}
\end{figure}

\secn{2.3 贸易模型回答的问题与关键均衡条件}

\textbf{贸易模型回答的四类问题（第 10 页）}：
\begin{enumerate}
  \item \term{贸易为什么发生？（Why trade happens，老师圈画强调）}——自给自足下\term{相对价格不同}（differential relative prices under autarky）；
  \item \term{什么决定贸易模式（trade patterns，圈画）}——哪国出口哪种商品？（等价于：什么决定自给自足相对价格的差异？）答案候选：\term{禀赋/要素的丰裕度、生产技术、偏好}等；
  \item \term{贸易的影响是什么？}——商品价格、生产模式、要素价格等；
  \item \term{贸易的福利效应（相对自给自足）？}——预算线扩张、效用变化、实际收入变化等。
\end{enumerate}

\begin{thmbox}[关键条件：相对价格均等化（Equalization of Relative Prices，第 11--17 页）]
\textbf{数值例子}：两种商品 A、B；两国 H（home）与 F（foreign）；
$P_A^H = 2,\; P_B^H = 1,\; P_A^F = 3,\; P_B^F = 2$（板书作答：相对价格
$P_A^H/P_B^H = 2$，$P_A^F/P_B^F = 3/2$；老师手写``Rel.\ price of A is higher at H''）。

\textbf{相对价格的含义}：在 H 国，A 是 B 的 \kw{2 倍贵}；老师手写批注：``\kw{Opportunity cost} of 1 A is 2 B''——在 H 卖 1 单位 A 可买 2 单位 B，卖 1 单位 B 可买 $1/2$ 单位 A；在 F 卖 1 单位 A 可买 $3/2$ 单位 B，卖 1 单位 B 可买 $2/3$ 单位 A。\kw{相对价格描述的是市场上的交换比率}（rate of exchange on the market，荧光高亮）。
\end{thmbox}

\begin{exbox}[套利练习：``过手贸易''如何赚钱（第 14--16 页）]
设你是唯一能在两国间无成本运输的人（老师用游戏``传送''作类比）。$P_A^H/P_B^H = 2 > P_A^F/P_B^F = 3/2$：
\begin{enumerate}
  \item 在 H 卖 \kw{1} 单位 A，按 H 的相对价格换得 \kw{2} 单位 B；（老师手写：\term{出口到相对价格低处}…… 原文``Export: rel.\ price lower''）
  \item 把 \kw{2} 单位 B 运到 F；
  \item 在 F 卖 \kw{2} 单位 B，换得 $\frac{4}{3}$ 单位 A；（手写：\term{从相对价格高处进口}``Import: rel.\ price higher''；并手写 $1B \Rightarrow \frac{2}{3}A$）
  \item 把 $\frac{4}{3}$ 单位 A 运回 H——净赚 $\frac{1}{3}$ 单位 A 的利润。
\end{enumerate}
老师手写：\kw{套利（Arbitrage）}、``Roundabout Trade''（迂回贸易）。手写标注两处关键动作（圈+叉）：``卖 A 于 H''（出口）与``运 A 回 H''（进口）是盈亏的两个环节。

\textbf{套利如何推动价格均等化（第 16 页）}：套利者不断``在相对价格高处卖、低处买''：
在 H 卖 A 买 B $\Rightarrow$ \kw{H 的 A 相对价格 $\downarrow$}（荧光高亮）；在 F 卖 B 买 A $\Rightarrow$ \kw{F 的 A 相对价格 $\uparrow$}（荧光高亮）。套利持续到两国 A 的相对价格相等为止。老师手写：``\kw{Equilibrium condition: rel. price equalize}''（均衡条件：相对价格均等化）。
\end{exbox}

\begin{intubox}[重要讨论（第 17 页）]
\term{相对价格对汇率是中性的（neutral to exchange rate，荧光高亮）}！——相对价格只取决于两种商品的交换比率，与用哪种货币（或汇率）计价无关。
\end{intubox}

\secn{2.4 禀赋模型：数值例子（自给自足）}

\begin{assumpbox}[禀赋模型设定（第 19--21 页）]
\begin{itemize}
  \item 两种商品：小麦（wheat, $W$）与布（cloth, $C$）；\kw{无生产}（老师手写``No production''于标题下）；
  \item 禀赋：$Q_w = 100$，$Q_c = 200$ $\Rightarrow$ PPF 退化为一个\term{点} $(100, 200)$（板书见图~\ref{fig:ch02-endppf}：PPF 是一个点，禀赋模型中固定）；
  \item 消费 $C_w, C_C$，价格 $P_w, P_C$；
  \item 自给自足下：$C_w = Q_w$，$C_c = Q_c$（消费点=禀赋点）；
  \item \kw{Cobb--Douglas 偏好（荧光高亮）}：$U = C_w^{0.2} C_C^{0.8}$。
\end{itemize}
\end{assumpbox}

\begin{figure}[H]
  \centering
  \includegraphics[width=0.55\linewidth]{figs/ch02_fig_05.png}
  \caption{第 19 页板书：禀赋经济中``PPF''退化为禀赋点 $(Q_w, Q_c) = (100, 200)$。}
  \label{fig:ch02-endppf}
\end{figure}

\begin{intubox}[无差异曲线与全微分（第 20 页手写）]
同一条无差异曲线上效用水平相同：$U(C_w, C_C) = \bar{U}$。全微分：
$MU_w\, dC_w = -MU_C\, dC_C \;\Rightarrow\; \dfrac{dC_C}{dC_w} = -\dfrac{MU_w}{MU_C}$
——无差异曲线斜率 = 边际效用之比的负值。
\end{intubox}

\begin{thmbox}[自给自足均衡的``三步解法''（第 21 页手写，老师称为``Dealer''）]
收入来自把全部禀赋卖给``交易商''（``Dealer''，手写），因此预算平衡要求 $I = P_w C_w + P_c C_c$。手写三步：
\begin{enumerate}
  \item \term{把全部产品卖给 Dealer}：$I = P_w Q_w + P_c Q_c$；
  \item \term{从 Dealer 处购买}：$\max_{C_w, C_C} U = C_w^{0.2} C_C^{0.8}$ s.t.\ $P_w C_w + P_c C_c = I$；
  \item \term{自给自足均衡条件}：$C_w = Q_w$，$C_C = Q_C$（生产决策点=消费决策点）。
\end{enumerate}
\end{thmbox}

\begin{figure}[H]
  \centering
  \includegraphics[width=0.5\linewidth]{figs/ch02_fig_06.png}
  \caption{第 21 页手写：禀赋模型自给自足的``三步解法''（卖清禀赋 $\to$ 预算内最优购买 $\to$ 消费=禀赋）。}
  \label{fig:ch02-dealer}
\end{figure}

\textbf{最优预算线：经过禀赋点（PPF），并与无差异曲线相切（第 22 页）}——老师手写``相切''条件：无差异曲线斜率 $-\frac{MU_1}{MU_2} = -\frac{P_1}{P_2}$（预算线斜率），即 $MU_1/P_1 = MU_2/P_2$（一手写注：两边乘 $1/P_i$ 后相等）。

\textbf{一阶条件与 C--D 支出法则（第 23 页）}：
\[
\frac{MU_C}{P_c} = \frac{MU_w}{P_w}
\qquad\Longrightarrow\qquad
{\color{hlorange}P_w C_w = 0.2 I},\quad {\color{hlorange}P_c C_c = 0.8 I}
\]
（手写将三式联立：$P_w C_w = 0.2I$，$P_c C_c = 0.8I$，$P_w Q_w + P_c Q_C = I$，并利用 $Q_w = C_w,\ Q_C = C_C$。）

\begin{warnbox}[PPT 红字提示（第 23 页）]
对 C--D 效用 $U = \prod_{g=1}^{G} C_g^{\alpha_g}$，永远有
\[
{\color{hlorange}C_g P_g = \frac{\alpha_g}{\sum_{i=1}^{G} \alpha_i}\, I}
\]
（荧光高亮）：\term{商品 $g$ 上的支出与 $\alpha_g$ 成正比}。
\end{warnbox}

\begin{figure}[H]
  \centering
  \includegraphics[width=0.62\linewidth]{figs/ch02_fig_07.png}
  \caption{第 24 页整页手写：C--D 最优消费的完整证明——由 FOC $\frac{MU_1}{P_1} = \cdots = \frac{MU_g}{P_g}$，令其等于 $\frac{1}{\lambda}$，得 $C_g P_g = \alpha_g \lambda$；代入预算约束 $\sum_g C_g P_g = I$ 解出 $\lambda = I / \sum_g \alpha_g$，故 $C_g P_g = \frac{\alpha_g}{\sum_g \alpha_g} I$：商品 $g$ 的支出占总收入的比例 = 其在 C--D 效用中的幂次 $\alpha_g$。}
  \label{fig:ch02-cdproof}
\end{figure}

\textbf{解自给自足价格（第 25--27 页）}：解得 $P_w = 0.2I/100$，$P_c = 0.8I/200$，$I = 100 P_w + 200 P_c$。老师手写指出：\term{2 个独立方程、3 个未知数}——只能解出相对价格：
\[
\frac{P_w}{P_C} = \frac{0.2/100}{0.8/200} = {\color{hlorange}\frac{1}{2}}.
\]
市场通过\kw{相对价格的调整}（荧光高亮）达到均衡，而不是（多个）绝对价格！可以把所有价格与收入同乘任意正数仍是均衡——回忆中级微观一般均衡中的\term{瓦尔拉斯定理}。我们只有 2 个市场、1 个``有效''价格：\kw{把另一种价格当作计价物（numeraire，荧光高亮）}（如令 $P_c = 1$）。板书（第 27 页）：令 $P_c = 1$ 解得 $P_w = \frac{1}{2}$，$P_c = 1$，$I = 250$——这解释了为什么贸易理论中通常使用\term{相对价格}。

\secn{2.5 两国贸易均衡与世界价格}

\begin{assumpbox}[两国设定（第 29--31 页）]
两国偏好相同（$U = C_w^{0.2} C_C^{0.8}$）；Home 禀赋 $Q_w^H = 100$，$Q_C^H = 200$；Foreign 禀赋 $Q_w^{Fa} = 200$，$Q_C^{Fa} = 200$。\redkw{F 相对更丰裕于小麦（wheat）}。

\textbf{自给自足价格}：$P_w = 0.2 I / Q_w$，$P_c = 0.8 I / Q_c$（以布为计价物 $P_c = 1$）：
\begin{itemize}
  \item H 国：$P_w^{Ha} = 0.5$，$P_c^{Ha} = 1$（手写标注``autarky''）；
  \item F 国：$P_w^{Fa} = 0.25$，$P_c^{Fa} = 1$；手写验算：$\left(\frac{P_w}{P_c}\right)^{Fa} = \frac{0.2 I / 200}{0.8 I / 200} = \frac{1}{4}$。
\end{itemize}
小麦在 \kw{H} 相对更贵（相对 F）——自给自足下 \kw{F} 出口小麦到 \kw{H}；反之 \kw{H} 出口布到 \kw{F}。\redkw{贸易的原因：两国禀赋的差异}（老师红笔下划线）。
\end{assumpbox}

\begin{thmbox}[贸易模式命题（第 31 页，荧光高亮原文）]
{\color{hlorange}在某种商品的禀赋上（相对）更丰裕的国家，把该商品出口给在该商品禀赋上相对更稀缺的国家。}
\end{thmbox}

\textbf{世界贸易均衡（第 32--34 页）}：把世界当作一个更大的经济体：世界偏好仍为 $U = C_w^{0.2} C_C^{0.8}$（$C_w = C_w^H + C_w^F$，$C_c = C_c^H + C_c^F$）；世界总禀赋
\[
Q_w = Q_w^H + Q_w^F = 300, \qquad Q_C = Q_C^H + Q_C^F = 400.
\]
两国价格必须均等（无套利）：$P_w = P_w^H = P_w^F$，$P_c = P_c^H = P_c^F$；仍需国际计价物，继续令 $P_c = 1$（现实中通常由美元扮演该角色）。由 C--D 支出法则（世界层面）：
\[
P_w = \frac{0.2 I}{300}, \quad P_c = \frac{0.8 I}{400}
\;\Longrightarrow\;
\frac{P_w}{P_c} = {\color{hlorange}\frac{1}{3}}
\]
（板书手写）。手写不等式（关键结论）：
\[
\underbrace{\left(\tfrac{P_w}{P_c}\right)^{\text{autarky, 出口国}}}_{=1/4}
\;\le\;
\underbrace{\left(\tfrac{P_w}{P_c}\right)^{\text{trade}}}_{=1/3}
\;\le\;
\underbrace{\left(\tfrac{P_w}{P_c}\right)^{\text{autarky, 进口国}}}_{=1/2}
\]
世界相对价格 $1/3$ 介于 F 的自给自足相对价格 $0.25$ 与 H 的 $0.5$ 之间。

\begin{thmbox}[一般化命题（第 34 页，荧光高亮）]
{\color{hlorange}均衡中的世界价格总是（弱地）介于两国的自给自足相对价格之间。}
\end{thmbox}

\textbf{Home 的贸易均衡（第 35--39 页）}：贸易下新预算线为
\[
P_w C_w^H + P_c C_c^H = P_w Q_w^H + P_c Q_c^H
\quad(\text{板书：}\; I = 200 P_c + 100 P_w = P_c \cdot C_c + P_w \cdot C_w).
\]
贸易后小麦变得相对更便宜；\term{预算线扩张表明福利改善}（老师手写``revealed preference''（显示偏好）指向该句，并手写两条理由：\kw{① 效用水平提高；② 选择范围扩张（原来的消费仍可保留）}）。

\begin{exbox}[``Gain from Trade''（第 36 页，荧光高亮）]
贸易得益源于：小麦的稀缺通过与一个在该商品上更丰裕的伙伴贸易得到补偿。图~\ref{fig:ch02-gain}：自给自足相对价格（预算线）为 $0.5$，贸易相对价格为 $1/3$；贸易下消费点从 PPF 上的禀赋点移到更高无差异曲线上（出口布、进口小麦）。
\end{exbox}

\begin{figure}[H]
  \centering
  \includegraphics[width=0.62\linewidth]{figs/ch02_fig_08.png}
  \caption{第 36 页：自给自足 vs.\ 贸易（Home）——PPF（禀赋点）、自给自足价格线（小麦相对价格 0.5）与贸易价格线（小麦相对价格 $1/3$），贸易线更平（小麦更便宜）；板书标注``$\downarrow$ export（出口布）''与``$\to$ import（进口小麦）''。}
  \label{fig:ch02-gain}
\end{figure}

\textbf{求解消费与贸易流（第 37--38 页）}：均衡时 $P_w = 1/3$，$P_c = 1$，$Q_w^H = 100$，$Q_C^H = 200$。收入 $I = P_w Q_w^H + P_c Q_c^H = \frac{1}{3} \times 100 + 200 = \frac{700}{3}$（板书验算）。由 $P_w = 0.2 I / C_w^H$：
\[
\frac{1}{3} = 0.2 \times \frac{700/3}{C_w^H} \;\Longrightarrow\; {\color{hlorange}C_w^H = 140}, \qquad {\color{hlorange}C_C^H = 186.7}.
\]
于是
\[
\text{布的出口 } EX_c = Q_c^H - C_c^H = {\color{hlorange}13.3}, \qquad
\text{小麦的进口 } IM_w = C_w^H - Q_w^H = {\color{hlorange}40}.
\]

\begin{figure}[H]
  \centering
  \includegraphics[width=0.6\linewidth]{figs/ch02_fig_09.png}
  \caption{第 38 页：贸易均衡下的 Home——消费点 $(140, 186.7)$ 在贸易预算线上、禀赋（生产）点 $(100, 200)$ 在 PPF 上；出口布 13.3、进口小麦 40。}
  \label{fig:ch02-consumption}
\end{figure}

\begin{thmbox}[贸易平衡（Trade Balance，第 39 页）]
出口价值 $EV_c = P_c (Q_c^H - C_c^H) = 1 \times 13.3 = 13.3$；
进口价值 $IV_w = P_w (C_w^H - Q_w^H) = \frac{1}{3} \times 40 = 13.3$（板书验算）。
两者相抵：\term{贸易平衡（trade balance）}。手写推导：预算约束 $P_c Q_c + P_w Q_w = P_c C_c + P_w C_w$ 移项得
\[
\underbrace{P_c (Q_c - C_c)}_{EV_c} - \underbrace{P_w (C_w - Q_w)}_{IV_w} = 0 .
\]
\end{thmbox}

\begin{sumbox}[禀赋模型小结（第 40 页原文红字）]
\begin{itemize}
  \item 各国\redkw{自给自足相对价格}不同，原因是\redkw{相对禀赋丰裕度}不同；
  \item \redkw{贸易模式}：出口自给自足下相对更便宜的商品（即在禀赋上相对更丰裕），进口禀赋相对更稀缺的商品；
  \item \term{贸易为什么能发生？}\redkw{贸易商靠套利赚钱！}
  \item \term{国家为什么允许贸易？}\redkw{两国福利都从贸易中得益！}
\end{itemize}
\end{sumbox}

\secn{2.6 一般化：$G$ 种商品、$N$ 个国家}

\begin{assumpbox}[一般化设定（第 42--43 页）]
$G$ 种商品（$g = 1, \dots, G$；手写补充国家 $i = 1, \dots, N$）；每个经济体对每种商品有禀赋 $Q_g^i$；各国代表性主体偏好相同且为 Cobb--Douglas：
\[
U^i = \prod_{g=1}^{G} \bigl( C_g^i \bigr)^{\beta_g}, \qquad \sum_{g=1}^{G} \beta_g = 1,
\qquad\text{支出份额：}\; P_g^i C_g^i = \beta_g I^i .
\]
\end{assumpbox}

\textbf{世界经济（第 44--45 页）}：所有国家开放（无成本）贸易后``世界经济''中\term{（相对）价格处处相同}。由 C--D 偏好（且各国偏好相同）得世界价格
\[
{\color{hlorange}P_g^w = \beta_g I / Q_g^w},
\]
其中 $P_g^w$ 是\kw{世界经济中的价格}（荧光高亮）、$I$ 是\term{世界收入}、$Q_g^w$ 是\kw{总禀赋（各国禀赋之和）}（荧光高亮）。

贸易下 Home 的收入为 $I^H = \sum_g P_g^w\, Q_g^H$（板书）。

\textbf{出口还是进口？（第 46--48 页）}：消费小于禀赋则出口，大于则进口。$g$ 的出口价值（$EV_g^H < 0$ 意味着进口）：
\[
EV_g^H = P_g Q_g^H - P_g C_g^H = P_g Q_g^H - \beta_g I^H
\]
（手写代入 $P_g C_g^H = \beta_g I^H$ 与 $P_g = \beta_g I / Q_g^w$）。用代数技巧 $I^H = I \cdot \frac{I^H}{I}$，两边除以 $I$：
\[
\frac{EV_g^H}{I} = \beta_g \left( \frac{Q_g^H}{Q_g^w} - \frac{I^H}{I} \right)
\quad\Longrightarrow\quad
\boxed{\;EV_g^H = \beta_g I \left( \frac{Q_g^H}{Q_g^w} - \frac{I^H}{I} \right)\;}
\]

\begin{thmbox}[核心结论（第 48 页，红字原文）]
一国出口商品 $g$ 当且仅当它在该商品的\redkw{世界禀赋中所占份额}高于其\redkw{世界收入中所占份额}：
\[
\frac{Q_g^H}{Q_g^w} > \frac{I^H}{I}.
\]
（$+$ 为出口，$-$ 为进口——手写笔记页第 59 页补充的等价写法：
$EX_i^g = Q_i^g - C_i^g = Q_i^g - \beta_g I^i / P_g^w$，
$EV_i^g = P_i^w EX_i^g = \beta_g I \left( \frac{Q_i^g}{Q_g^w} - \frac{I^i}{I} \right)$。）
\end{thmbox}

\begin{intubox}[贸易平衡的一般证明（第 49 页手写）]
$\displaystyle \sum_{g=1}^G EV_g^H = 0$。证明：$\sum_g EV_g^H = \underbrace{\sum_g P_g Q_g^H}_{=\,I^H\ (\text{定义})} - \underbrace{\sum_g \beta_g}_{=1\ (\text{第 43 页})} I^H = I^H - I^H = 0$。
手写注：这仍是\term{国家预算约束}的等价形式。
\end{intubox}

\textbf{应用：石油贸易（第 50--53 页）}。关键表达式 $EV_g^H = \beta_g I \left( \frac{Q_g^H}{Q_g^w} - \frac{I^H}{I} \right)$：
\begin{itemize}
  \item 一国拥有世界 \kw{10\%} 的石油，但只有世界 \kw{1\%} 的收入 $\Rightarrow$ \kw{出口}价值等于世界石油消费 \kw{9\%} 的石油：$\frac{EV_g^H}{\beta_g I} = 10\% - 1\% = 9\%$（板书作答）；
  \item 一国拥有世界 10\% 的石油，但有世界 20\% 的收入 $\Rightarrow$ \kw{进口}价值等于世界石油消费 10\% 的石油：$10\% - 20\% = -10\%$（板书作答）。
\end{itemize}

\begin{table}[H]
  \centering
  \caption{石油贸易：模型估计 vs.\ 现实（2014 年相对产量与 GDP；净出口单位：千桶/日）。数据来源：EIA of the United States, The World Bank, IEA（PPT 第 53 页）；老师手写标注表头对应 $\frac{Q_g}{Q_g^w}$ 与 $\frac{I}{I^w}$，并注``净出口估计''列 $= \left( \frac{Q_g}{Q_g^w} - \frac{I}{I^w} \right) \times \beta_g I$ $=$（第 2 列 $-$ 第 3 列）$\times$ 86.23 百万桶/日（2012 年世界产量）。}
  \label{tab:ch02-oil}
  \begin{tabular}{@{}lrrrr@{}}
    \toprule
    经济体 & 石油产量相对世界 & GDP 相对世界 & 净出口估计 & 2012 实际净出口 \\
    \midrule
    哥斯达黎加 & 0.0003\% & 0.0633\% & $-54.3$ & $-46.3$ \\
    法国 & 0.0656\% & 3.6142\% & $-3{,}060.0$ & $-1{,}943.2$ \\
    沙特阿拉伯 & 12.4966\% & 0.9533\% & $9{,}953.8$ & $10{,}546.0$ \\
    西班牙 & 0.0430\% & 1.7940\% & $-1{,}509.9$ & $-1{,}557.7$ \\
    \bottomrule
  \end{tabular}
\end{table}

模型估计的净出口方向与量级与实际数据相当接近——禀赋差异确实是石油贸易格局的主要解释因素。

\secn{2.7 模型构建方法地图（据课后手写笔记整理）}

\begin{mapbox}[禀赋模型构建方法地图（课后手写笔记第 54--61 页的框架）]
\textbf{1. 供给与需求假定 \& 最优化条件}
\begin{itemize}
  \item \term{供给端}：PPF、生产决策点（在 PPF 上）$(Q_1, Q_2)$——\kw{禀赋模型中固定}（PPF 退化为一个点）；
  \item \term{需求端}：无差异曲线（与效用函数），消费决策点 $(C_1, C_2)$ 由 \kw{IC 与预算线相切}决定；
  \item \term{预算线（``价格线''）}：$(B.C.)\ I = P_1 C_1 + P_2 C_2$，\kw{经过生产决策点}——左边 $I = P_1 Q_1 + P_2 Q_2$ 的含义：\term{生产出的产品出售形成收入}；
  \item 变形：$P_1 (C_1 - Q_1) + P_2 (C_2 - Q_2) = 0$，含义：① 几何上预算线经过点 $(Q_1, Q_2)$；② \term{购买一种产品的支出来自卖出另一种产品的收入}；
  \item 特例（C--D 效用 $U = C_1^{\alpha_1} C_2^{\alpha_2}$，$\alpha_1 + \alpha_2 = 1$）：$\frac{C_1 P_1}{I} = \alpha_1$，$\frac{C_2 P_2}{I} = \alpha_2$——\term{产品 $i$ 的开支占总收入的比例 = 其在效用函数中的幂次 $\alpha_i$}。
\end{itemize}
\textbf{2. 自给自足（Autarky）与贸易均衡}
\begin{itemize}
  \item \term{2.1 自给自足}：条件是\kw{生产点与消费点重合}——PPF 与价格线相切决定（禀赋模型中进一步简化为固定点）；价格线与无差异曲线相切决定（$\frac{MU_1}{MU_2} = \frac{P_1}{P_2}$）。求解：自给自足相对价格 $P_1/P_2$（或以商品 2 为 numeraire，令 $P_2 = 1$ 求 $P_1$）。特例：$U = C_1^{\alpha_1} C_2^{\alpha_2}$ \& $C_1 = Q_1, C_2 = Q_2$ $\Rightarrow$ $P_1 = \frac{\alpha_1 I}{Q_1}$，$P_2 = \frac{\alpha_2 I}{Q_2}$；
  \item \term{2.2 贸易下均衡（一国情形）}（参见课件第 10--14 页）：均衡条件 ① \kw{无套利：相对价格均等}；② \term{出口 $=$ 生产 $-$ 消费，进口 $=$ 消费 $-$ 生产}——此时 $(C_1, C_2)$ 与 $(Q_1, Q_2)$ 不必重合，但都在预算线上。任何 $P_1/P_2$ 都可解出 $C_1, C_2, Q_1, Q_2$ 为 $P_1/P_2$ 的函数；进一步 $Q_1 - C_1$ 与 $Q_2 - C_2$ 亦为其函数：若 $Q_i - C_i > 0$ 出口，$< 0$ 进口；
  \item \term{2.3 贸易模式（Patterns of Trade）}：取决于自给自足下两国相对价格——每个国家\term{出口相对价格低的产品、进口相对价格高的产品}。Endowment 下相对价格由禀赋相对丰裕度（$Q_1/Q_2$）决定：\kw{出口禀赋相对（另一国）更丰裕的产品}；
  \item \term{2.4 贸易下的均衡（世界市场出清）}：一般解法——记商品 1 为 H 的出口品（F 的进口品），对 H、F 两国重复 2.2：Export Supply $= Q_1^H - C_1^H$ 是 $P_1/P_2$ 的函数，Import Demand $= C_1^F - Q_1^F$ 是 $P_1/P_2$ 的函数；\kw{令二者相等}（一国的出口即另一国的进口）即可解出 $P_1/P_2$。注：由 \term{Walras 定理}只需考察一种商品。Endowment 特例：把全球看作一个经济体 $P_g = \frac{\alpha_g I^w}{Q_g^w}$（$g = 1, 2$），$I^w$ 为全球总收入、$Q_g^w$ 为全球总禀赋；
  \item \term{2.5 Balance of Trade（任一国家）}：出口额 $EV = P_1 (Q_1 - C_1)$，进口额 $IV = P_2 (C_2 - Q_2)$；$EV = IV$ 即等价于预算线方程——\term{只要预算中的收入来自卖出产品，贸易模型总推出贸易平衡}（无顺差也无逆差）；
  \item \term{2.6 Gains from Trade}（以 H 出口 1 进口 2 为例）：贸易后预算线绕禀赋（生产）点向外转动（一般情形相对价格比自给自足``增加''——原因：两国相对价格趋于相等），消费点移到出口品更少、进口品更多的位置，\kw{效用上升}；原来的自给自足消费束仍可负担（选择集扩大）。Endowment 下的图示：贸易使效用上升。
\end{itemize}
\textbf{3. Generalized Endowment Model}：$g = 1, \dots, G$ 种产品、$i = 1, \dots, N$ 个国家；$Q_i^g$ 禀赋、$I^i = \sum_{g=1}^G P_i^g Q_i^g$ 收入、效用 $U = \prod_g (C_i^g)^{\beta_g}$（$\sum \beta_g = 1$）。
\begin{itemize}
  \item \term{3.1 Autarky}：$C_i^g = Q_i^g$；联立 C--D 支出式 $P_i^g C_i^g = P_g I^i$ $\Rightarrow$ $P_i^g = \frac{P_g I^i}{Q_i^g}$（Autarky Price）；
  \item \term{3.2 Trade}：世界作为大经济体：$I = \sum_i I^i$，$Q_g^w = \sum_i Q_i^g$；利用 C--D 形式得 $P_i^w = \frac{\beta_g I^w}{Q_g^w}$（Trade Equilibrium Price，贸易均衡价格）；
  \item \term{3.3 Trade Patterns}：$EX_i^g = Q_i^g - C_i^g = Q_i^g - \frac{\beta_g I^i}{P_i^w}$；出口额 $EV_i^g = P_i^w EX_i^g = \beta_g I \left( \frac{Q_i^g}{Q_g^w} - \frac{I^i}{I} \right)$（$+$ 为出口，$-$ 为进口）。
\end{itemize}
\textbf{4. Edgeworth 盒视角}（笔记第 61 页）：把 H、F 两国画在 Edgeworth Box 的左下、右上两端——Endowment Model 实际上就是中级微观``一般均衡''章节分析的\term{交换经济}；``Gains from Trade for both Countries''（两国都从贸易得益）\kw{等价于福利经济学第一定理}。
\end{mapbox}

\begin{figure}[H]
  \centering
  \includegraphics[width=0.52\linewidth]{figs/ch02_fig_12.png}
  \caption{课后手写笔记框架页（一）：贸易模型结构（禀赋模型）——供给/需求端设定、预算线（价格线）的双重含义、C--D 特例。}
  \label{fig:ch02-note1}
\end{figure}

\begin{figure}[H]
  \centering
  \includegraphics[width=0.5\linewidth]{figs/ch02_fig_10.png}
  \caption{课后手写笔记：自给自足的图示——上：一般情形（PPF 与价格线、价格线与 IC 相切，生产点=消费点）；下：Endowment 模型情形（``PPF''退化为一点）。}
  \label{fig:ch02-note2}
\end{figure}

\begin{figure}[H]
  \centering
  \includegraphics[width=0.62\linewidth]{figs/ch02_fig_11.png}
  \caption{课后手写笔记：把禀赋模型画进 Edgeworth 盒——H、F 分别位于左下与右上，禀赋 $Q_1^H, Q_2^H$ 与 $Q_1^F, Q_2^F$ 张成整盒；贸易均衡价格线下两国的消费点通过``H 出口/F 进口''衔接；两国贸易得益 $\Leftrightarrow$ 福利经济学第一定理。}
  \label{fig:ch02-note3}
\end{figure}

\begin{sumbox}[本讲总结：通用结论与方法]
\begin{enumerate}
  \item \term{建模语法}：模型 $=$ 设定（谁在决策、选择集、交易条件）$+$ 目标与约束 $+$ 均衡概念；先分清内生/外生变量（厂商例子：$L_i, K_i, Q_i, P_i$ 内生；$w_i, r_i, F, P(Q_i)$ 外生）。
  \item \term{古典贸易模型共同设定}：$\ge 2$ 国（另一国可记为 ROW）、$\ge 2$ 种商品（一种商品无法平衡预算）；供给面 $=$ 禀赋 $+$ 技术 $\to$ PPF $+$ 完全竞争（MC=P，w=MRPL，r=MRPK）；需求面 $=$ 同质位似偏好下的代表性主体（IC $+$ 预算约束）。
  \item \term{关键均衡条件}：开放贸易后\textbf{相对价格均等化}（无套利）；套利是价格均等化的微观机制（低买高卖的``迂回贸易''）；相对价格 $=$ 市场交换比率 $=$ 机会成本，对汇率中性。
  \item \term{C--D 效用的工作马}：$U = \prod_g C_g^{\alpha_g}$ $\Rightarrow$ $C_g P_g = \frac{\alpha_g}{\sum_i \alpha_i} I$——支出份额等于幂次；这是所有解析求解的引擎（板书有完整证明）。
  \item \term{价格只定相对值}：$N$ 个市场只有 $N-1$ 个独立价格（瓦尔拉斯定理）；把一种商品记为 numeraire；均衡由相对价格调整实现。
  \item \term{禀赋模型三段论}：禀赋差异 $\to$ 自给自足相对价格差异 $\to$ 贸易模式（出口相对丰裕品）；世界均衡价格介于两国自给自足价格之间（出口国自给价 $\le$ 世界价 $\le$ 进口国自给价）。
  \item \term{贸易必然平衡}：预算约束 $\Rightarrow$ $\sum_g EV_g = 0$（无顺差逆差）；\term{两国都得益}（显示偏好：原消费束仍可负担 $+$ 选择集扩大）。
  \item \term{一般化核心公式}：$EV_g^i = \beta_g I \left( \frac{Q_g^i}{Q_g^w} - \frac{I^i}{I} \right)$——出口该商品 $\iff$ 该商品的世界禀赋份额 $>$ 世界收入份额；石油贸易数据验证良好。
\end{enumerate}
\end{sumbox}
